\chapter{RESUMO}\label{chap:resumo}

A agricultura moderna demanda mecanismos de acompanhamento capazes de apoiar a tomada de decisão em áreas cultivadas, especialmente quando a extensão da lavoura, a variabilidade climática e o custo de inspeções presenciais dificultam o monitoramento frequente. Nesse contexto, o presente trabalho propõe o desenvolvimento de um sistema de monitoramento agrícola baseado em imagens de satélite, com uso de sensoriamento remoto, índices espectrais e dados climáticos gratuitos. O objetivo é desenvolver um protótipo funcional capaz de selecionar uma área de interesse, consultar imagens Sentinel-2 por meio do Google Earth Engine, calcular índices como o \textit{Normalized Difference Vegetation Index} e o \textit{Normalized Difference Water Index}, gerar mapas temáticos, produzir séries temporais e indicar possíveis anomalias na vegetação. A metodologia prevê a implementação do sistema em Python, com interface em Streamlit, integração com geemap e Plotly, processamento geoespacial em nuvem pelo Google Earth Engine e consulta complementar a dados climáticos da NASA POWER. A validação será realizada por meio de testes funcionais, análise de três áreas agrícolas representativas, comparação dos valores dos índices com faixas esperadas na literatura e interpretação conjunta com variáveis climáticas, como precipitação, temperatura e radiação solar. Espera-se demonstrar a viabilidade técnica de uma solução acadêmica de baixo custo, reproduzível e baseada em ferramentas abertas para apoiar a análise da condição da vegetação em lavouras, sem substituir a avaliação agronômica em campo.
